\documentclass[11pt,a4paper]{article}

\usepackage[T1]{fontenc}
\usepackage[utf8]{inputenc}
\usepackage{geometry}
\usepackage{enumitem}
\usepackage{longtable}
\usepackage{array}
\usepackage{xcolor}
\usepackage{hyperref}

\geometry{margin=2.5cm}
\hypersetup{colorlinks=true, linkcolor=blue, urlcolor=blue}
\setlist[itemize]{leftmargin=*, itemsep=0.3em}
\setlist[enumerate]{leftmargin=*, itemsep=0.35em}
\newcommand{\loc}[1]{\texttt{#1}}
\newcommand{\sev}[1]{\textbf{#1}}

\title{LLM Copy-Paste Residue Check}
\author{Scan report for the active \texttt{thesis.tex} dependency tree}
\date{12 July 2026}

\begin{document}
\maketitle

\section*{Scope}

This report checks the active dissertation source for phrases that often appear
when text has been pasted from an LLM chat or from generic LLM-style rewrites.
The check covers the active \loc{thesis.tex} dependency tree only. Existing
thesis files were not modified.

This is not an AI-authorship detector. The goal is narrower and more useful:
find chat residue, prompt/answer residue, generic filler, tutorial tone, and
source scaffolding that should be removed before final submission.

\section*{Summary}

\begin{itemize}
  \item \sev{No direct smoking-gun chat leakage found.} I did not find active
  occurrences of phrases such as \verb|here is|, \verb|as requested|,
  \verb|according to your definition|, \verb|based on your definition|,
  \verb|using your definition|, \verb|your thesis|, or \verb|let me|.
  \item \sev{Some false-positive direct matches were harmless.} For example,
  \loc{chapterconflict.tex:529} and \loc{forwardtight.tex:21} use the word
  ``below'' in normal academic prose.
  \item \sev{Real issues exist in three clusters:} prompt/response residue in the
  introduction's LLM example, generic/tutorial phrasing in several technical
  sections, and source-scaffolding language around the framework-user discussion.
  \item \sev{Highest-priority cleanup:} review \loc{einleitung.tex:67--81},
  \loc{chapterconflict.tex:107}, \loc{chapterconflict.tex:1445},
  \loc{chapterconflict.tex:1502}, \loc{tightmotivation.tex:673--738},
  \loc{qviol.tex:868}, and \loc{fazit.tex:10--38}.
\end{itemize}

\section*{Pattern Counts}

After stripping fully commented lines, the scan found the following counts:

\begin{center}
\begin{tabular}{lr}
\textbf{Category} & \textbf{Matches} \\
\hline
Direct chat leakage & 2 \\
Prompt/response residue & 13 \\
Generic transitions & 24 \\
Vague filler & 17 \\
Hedging & 5 \\
Instructional/tutorial tone & 18 \\
Scaffolding or source debris & 2 \\
Citation/fact red flags & 23 \\
\end{tabular}
\end{center}

The direct-chat count is misleading: both matches were ordinary uses of
``below'', not actual chat leakage.

\section*{Priority 1: Inspect and Rewrite}

\begin{enumerate}
  \item \sev{Prompt/response wording in the introduction.}
  \loc{einleitung.tex:67} and \loc{einleitung.tex:81} discuss a prompt, a response,
  and a hallucination detector. This may be intentional because the passage is
  about LLM hallucinations, but it reads close to chat-analysis language.
  Suggested action: make it explicit that this is a documented case study or
  demonstration, not an embedded chat transcript. Avoid vague phrases such as
  ``The solution uses'' unless the solution is named.

  \item \sev{``The user'' phrasing in formal framework text.}
  \loc{chapterconflict.tex:107}, \loc{chapterconflict.tex:1445},
  \loc{chapterconflict.tex:1502}, and \loc{conflictmeat.tex:1180} refer to
  ``the user'' of the framework. This is not necessarily wrong, but it sounds like
  tool-documentation prose rather than dissertation prose.
  Suggested action: replace with ``the designer'', ``the analyst'', ``the regulator'',
  or ``the specification engineer'', depending on context.

  \item \sev{Broken framework-user paragraph.}
  \loc{chapterconflict.tex:1502} combines user/tool language, generic terms, and
  an already-known broken sentence: ``Unless all the other tools in the framework
  are.'' This is one of the strongest copy/paste-scaffolding signals in the scan.
  Suggested action: rewrite the paragraph from scratch as an academic workflow.

  \item \sev{Tutorial examples with ``The user inputs''.}
  \loc{tightmotivation.tex:673--738} repeatedly uses phrases such as
  ``The user inputs''. This is acceptable for interactive examples, but it reads
  more like a generated tutorial than a formal dissertation section.
  Suggested action: use ``the trace begins with'', ``the next symbol is'', or
  ``the input word is extended by''.

  \item \sev{Question-answer setup in quantitative blame.}
  \loc{qviol.tex:868} asks: ``does adding up the two blame counters always give
  back the same score''. This conversational question style is not a direct LLM
  leak, but it reads like explanatory chat prose.
  Suggested action: rewrite as a formal motivation for the lemma and align it
  with the actual inequality proved.

  \item \sev{Conclusion still has generic LLM-like phrasing.}
  \loc{fazit.tex:10}, \loc{fazit.tex:20}, and \loc{fazit.tex:38} contain phrases
  such as ``clear and practical'', ``more expressive and easier to use'', and
  ``comprehensive solution''. These are plausible but generic.
  Suggested action: replace with concrete claims about decidability, finite-state
  monitor synthesis, conflict density, and blame attribution.
\end{enumerate}

\section*{Priority 2: Citation and Fact-Claim Red Flags}

These are not LLM leaks, but they are common LLM-style overclaims unless they are
properly sourced or carefully qualified.

\begin{longtable}{>{\raggedright\arraybackslash}p{0.24\linewidth}
                  >{\raggedright\arraybackslash}p{0.36\linewidth}
                  >{\raggedright\arraybackslash}p{0.30\linewidth}}
\textbf{Location} & \textbf{Flagged wording} & \textbf{Suggested action} \\
\hline
\loc{einleitung.tex:29} & ``Studies show that in places with many regulations...'' & Add a citation or make the claim more cautious. \\
\loc{einleitung.tex:33} & ``in the literature'', ``widely cited'', ``roughly'' & The footnotes help; keep exact source framing and avoid sounding broader than the cited reports. \\
\loc{einleitung.tex:43} & ``According to Legal Complex... hundreds of millions'' & The figure source helps; consider adding access date and exact scope. \\
\loc{chapterconflict.tex:401} & ``well known procrastination effect'' & Already cited; keep if the cited works explicitly support this term. \\
\loc{chapterconflict.tex:660} & ``commonly used'' & Harmless, but can be shortened. \\
\loc{chapterconflict.tex:688--690} & ``according to our definition(s)'' & Not chat leakage, but repetitive; replace with direct formal reference. \\
\loc{conflictmeat.tex:252} & ``According to our definition'' & Fine technically; consider ``By Definition~...''. \\
\loc{qviol.tex:114,145} & ``According to the definition...'' & Fine technically; ``By Definition~...'' is cleaner. \\
\end{longtable}

\section*{Priority 3: Generic or Tutorial Tone}

These passages are probably harmless, but they are places where the prose may
sound generated or overly didactic.

\begin{itemize}
  \item \loc{chapterconflict.tex:696}: ``We saw that... We proceed now to show...''
  could be tightened to ``We now prove that...''.
  \item \loc{chapterconflict.tex:774}: ``We now turn to...'' is normal but repeated
  often. Use sparingly.
  \item \loc{conflictmeat.tex:958}: ``We now show... very natural normative
  intuitions'' sounds conversational. Prefer a direct theorem-introduction sentence.
  \item \loc{chaptermonitor.tex:11}: ``We begin by motivating...'' is normal, but
  the paragraph already has several guidepost phrases. Reduce navigation prose.
  \item \loc{chaptermonitor.tex:348}, \loc{chaptermonitor.tex:1196},
  \loc{tightmotivation.tex:425}, \loc{tightmotivation.tex:665}, and
  \loc{tightmotivation.tex:751}: repeated ``Crucially'' can sound LLM-like.
  Keep only where the claim is genuinely a turning point.
  \item \loc{forwardtight2.tex:200}: ``Let us consider instead...'' is fine in a
  proof/example, but use consistently with the rest of the thesis style.
  \item \loc{quantitativeviol.tex:463}: ``We now turn to...'' is acceptable but can
  be shortened.
\end{itemize}

\section*{Probably Harmless Matches}

\begin{itemize}
  \item \loc{chapterconflict.tex:529}: ``The satisfaction relation below...'' is
  ordinary exposition.
  \item \loc{forwardtight.tex:21}: ``semantic clauses below'' is ordinary exposition.
  \item \loc{preliminaries.tex:51}, \loc{preliminaries.tex:276}, and
  \loc{preliminaries.tex:424}: ``Consider the following examples'' is normal in
  mathematical writing.
  \item \loc{declaration.tex:13--14}: AI tools are explicitly declared. This is not
  leakage; it is required disclosure text.
  \item \loc{einleitung.tex:67--81}: ``prompt'' and ``response'' may be necessary
  because the passage discusses an LLM hallucination example. It still deserves
  stylistic cleanup.
\end{itemize}

\section*{Recommended Manual Search Patterns}

Before final submission, run quick searches for these exact phrases:

\begin{verbatim}
here is
here are
as requested
according to your definition
based on your definition
using your definition
your thesis
your framework
the user
the prompt
the response
step-by-step
it is important to note
this highlights
this underscores
serves as a
provides a robust
clear and practical
comprehensive solution
roughly speaking
we now show
we now turn to
let us consider
TODO
placeholder
citation needed
\end{verbatim}

\section*{Bottom Line}

The scan did not find obvious active chat-copy residue such as ``here is the
rewritten paragraph'' or ``according to your definition''. The dissertation does,
however, contain several passages that read like generated explanatory prose or
tool documentation. The highest-value cleanup is to replace ``the user'' with a
domain role, rewrite the broken framework workflow paragraph, reduce tutorial
phrasing in tight-semantics examples, and make broad claims more precise or cited.

\end{document}
